% Avsnitt 1.1-1.6, ur lectures/L01/appendix/a_combinational_logic.md (A.1-A.6).

\section{Digitala signaler och logiska grindar}
\label{c1:sec:gates}

En digital signal antar ett av två värden, \code{0} eller \code{1}, i stället för det
kontinuerliga intervall en analog signal använder. Allt som ligger tillräckligt nära \code{0}
eller \code{1} läses som exakt \code{0} eller \code{1}, så små mängder elektriskt brus ändrar
ingenting. Den brusimmuniteten är anledningen till att nästan all modern beräkning, från
mikrokontroller till FPGA:er, är digital.

En \textbf{logisk grind} har en eller flera enbitarsingångar, en enbitarsutgång och en fast
boolesk funktion mellan dem.

\textbf{Kombinatorik}, det här kapitlets ämne, producerar utgångar enbart utifrån de aktuella
ingångarna, utan minne. \textbf{Sekvensnät} (\chapref{c3:ch}) beror på de aktuella ingångarna
\emph{och} på tillstånd lagrat i vippor.

Varje grundläggande grind med två ingångar, plus den unära \code{NOT}:

\begin{narrowtable}{cc*{6}{c}}
  \thead{A} & \thead{B} & \thead{AND} & \thead{OR} & \thead{NAND} & \thead{NOR} & \thead{XOR} &
    \thead{XNOR} \\
  \midrule
  0 & 0 & 0 & 0 & 1 & 1 & 0 & 1 \\
  0 & 1 & 0 & 1 & 1 & 0 & 1 & 0 \\
  1 & 0 & 0 & 1 & 1 & 0 & 1 & 0 \\
  1 & 1 & 1 & 1 & 0 & 0 & 0 & 1 \\
\end{narrowtable}

\begin{narrowtable}{cc}
  \thead{A} & \thead{NOT A} \\
  \midrule
  0 & 1 \\
  1 & 0 \\
\end{narrowtable}

Värt att lägga på minnet direkt ur tabellerna:
\begin{itemize}
  \item \code{NAND}, \code{NOR} och \code{XNOR} är inverserna av \code{AND}, \code{OR} och
    \code{XOR}.
  \item \code{XOR} är \code{1} när ingångarna skiljer sig åt. Bortom två ingångar generaliseras
    det till \textbf{udda paritet}: \code{1} när ett udda antal ingångar är \code{1}
    (Övning~\ref{c1:ex:xor3}).
  \item Varje boolesk funktion kan byggas av \code{AND}, \code{OR} och \code{NOT}, och även av
    enbart \code{NAND} eller enbart \code{NOR}.
\end{itemize}

% ------------------------------------------------------------------------------------------
\section{Boolesk algebra som notation}
\label{c1:sec:boolean}

Tre operatorer, som avbildas rakt av på grindarna ovan:

\begin{booktable}{@{}p{5em}L@{}}
  \thead{Grind} & \thead{Notation i boolesk algebra} \\
  \midrule
  AND & multiplikation: $A$ \textbf{AND} $B$ skrivs $AB$ \\
  OR & addition: $A$ \textbf{OR} $B$ skrivs $A + B$ \\
  NOT & prim: \textbf{NOT} $A$ skrivs $A'$ \\
\end{booktable}

En funktion skriven som en OR av AND-termer, till exempel $X = AB' + CD$, står på
\textbf{summa-av-produkter}-form.

Varje sanningstabell kan läsas av direkt som en summa av produkter:
\begin{enumerate}
  \item Leta upp varje rad där utgången är \code{1}.
  \item Skriv för var och en en AND-term som innehåller samtliga ingångar: direkt om ingången är
    \code{1} på den raden, primmad om den är \code{0}.
  \item OR:a ihop de termerna.
\end{enumerate}

Till exempel:

\begin{narrowtable}{ccc}
  \thead{A} & \thead{B} & \thead{X} \\
  \midrule
  0 & 0 & 0 \\
  0 & 1 & 1 \\
  1 & 0 & 1 \\
  1 & 1 & 0 \\
\end{narrowtable}

Raden $AB = 01$ bidrar med $A'B$ och raden $AB = 10$ med $AB'$, vilket ger $X = A'B + AB'$. Stäm
av mot tabellen: det är \code{XOR}, så samma nät är en enda \code{XOR}-grind i stället för två
\code{AND}-grindar, två \code{NOT}-grindar och en \code{OR}.

\subsection*{Lagarna, som en påminnelse}
Du har mött dem förut. De står här för att boken ska ha en enda notation för dem och för att
senare kapitel ska kunna hänvisa till dem, inte för att de behöver läras ut:

\begin{booktable}{@{}p{8em}p{9em}L@{}}
  \thead{Lag} & \thead{Med OR} & \thead{Med AND} \\
  \midrule
  Identitet & $A + 0 = A$ & $A \cdot 1 = A$ \\
  Dominans & $A + 1 = 1$ & $A \cdot 0 = 0$ \\
  Idempotens & $A + A = A$ & $A \cdot A = A$ \\
  Komplement & $A + A' = 1$ & $A \cdot A' = 0$ \\
  Absorption & $A + AB = A$ & $A(A + B) = A$ \\
  De Morgan & $(A + B)' = A'B'$ & $(AB)' = A' + B'$ \\
\end{booktable}

\textbf{De Morgan är den lag boken lutar sig mot}, mindre för minimering än som en regel för hur
en inversion flyttas genom en grind: en \code{OR} med båda ingångarna inverterade \emph{är} en
\code{NAND}, och en \code{AND} med båda ingångarna inverterade \emph{är} en \code{NOR}. Det är
därför en \code{NAND}-grind kan bygga vad som helst, vilket Övning~\ref{c1:ex:nand} låter dig
skriva i VHDL, och det kommer tillbaka varje gång en aktivt låg signal möter logik skriven för
aktivt hög. I en kurs där varje reset och varje knapp är aktivt låg är det ofta.

Glappet mellan en rakt avläst summa av produkter och det minsta nätet är vad ett
\textbf{Karnaughdiagram} sluter, och \secref{c2:sec:kmap} går igenom ett från början till slut.
Inte för att tekniken är ny för dig, utan för att det är den minimerade ekvationen som skrivs i
VHDL, och \chapref{c8:ch} härleder sin tillståndsmaskins nästa-tillstånd-logik på precis det
sättet.

% ------------------------------------------------------------------------------------------
\section{Att bygga och simulera en krets i CircuitVerse}
\label{c1:sec:cv}

CircuitVerse (\url{https://circuitverse.org/simulator}) är en gratis, webbläsarbaserad
logiksimulator. Varje grindnät i den här boken byggs och testas där för hand innan det skrivs i
VHDL:
\begin{enumerate}
  \item Öppna simulatorn och starta ett nytt projekt.
  \item Placera ett \textbf{input}-element per insignal och ett \textbf{output}-element per
    utsignal, omdöpta så att de matchar dina variabler.
  \item Placera de grindar din ekvation kräver, från panelen \emph{Gates}.
  \item Koppla ihop nätet en grind i taget, enligt din härledda ekvation. Du ska redan veta exakt
    vilka grindar du behöver innan du öppnar simulatorn, vilket är därför \secref{c1:sec:boolean}
    kommer först.
  \item Slå om varje ingång och kontrollera att utgången stämmer med din sanningstabell på
    \textbf{varje} rad.
  \item Spara via \code{Project -> Save Online}, eller \code{Project -> Download as image/data}
    för en återimporterbar \file{.cv}-fil.
\end{enumerate}

\textbf{Steg 5 är inte valfritt, och ingenting annat gör det åt dig.} Tre regler gör det till en
kontroll snarare än en formalitet:
\begin{itemize}
  \item \textbf{Uttömmande, inte stickprov.} Två ingångar är fyra rader, tre är åtta. Vid de
    storlekarna finns ingen ursäkt för att sampla. För sekvensnäten längre fram blir det här
    ”stega klockan och kontrollera vid varje flank”, samma disciplin tillämpad på tid.
  \item \textbf{Härled först, simulera sedan.} Rita kretsen och bestäm \emph{sedan} vad den ska
    göra, så kommer du att läsa av dina egna förväntningar ur simuleringen, vilket inte bevisar
    något.
  \item \textbf{En avvikelse är information.} Den säger att ritningen och ekvationen är oense, och
    att en av dem har fel.
\end{itemize}

Varje hårdvarubugg du kommer att jaga hittas genom att förutsäga ett beteende och sedan
kontrollera det. Att göra det för hand på en krets med fyra grindar är hur du lär dig göra det på
en du inte kan överblicka på en gång.

När konstruktionen väl är VHDL tar en självkontrollerande testbänk över det jobbet. Varje
genomarbetat exempel och varje övning har en utdelad; du kör dem snarare än skriver dem, vilket
\secref{c2:sec:testbenches} går igenom.

% ------------------------------------------------------------------------------------------
\section{Precis så mycket VHDL att du kan läsa och skriva ett grindnät}
\label{c1:sec:vhdl}

VHDL beskriver hårdvara, inte en sekvens av instruktioner. Varje sats nedan körs konkurrent och
kontinuerligt, precis som grindarna den beskriver. Språket kommer allteftersom det behövs; det här
avsnittet räcker för ett enkelt kombinatoriskt nät.

Varje VHDL-modul har två delar:
\begin{itemize}
  \item En \textbf{entitet}, vyn utifrån: portarna, och ingenting om vad som händer inuti.
  \item En \textbf{arkitektur}, implementationen bakom den.
\end{itemize}

Den uppdelningen är ett kontrakt, inte ceremoni. Entiteten är allt en annan konstruktion behöver
för att kunna använda modulen, vilket är det som låter en konstruktion instansiera en annan utan
att läsa den, och som låter dig skriva om en implementation utan att röra det som beror på den.
Formen är en header vid sidan av en källfil, eller ett gränssnitt vid sidan av sin klass;
skillnaden är att VHDL kräver den för varje modul, utan något sätt att hoppa över gränssnittet
eller läcka implementationen genom det.

Ta \code{or_gate}: ingångarna \code{a}, \code{b}, utgången \code{x}, och $x = a + b$.

\begin{vhdlcode}
entity or_gate is
    port(a, b: in std_logic;
         x   : out std_logic);
end entity;
\end{vhdlcode}

\lecfig[0.62]{lectures/L01/appendix/images/or_gate_entity.png}{Entiteten \code{or_gate}: vyn
utifrån, och ingenting om vad som händer inuti.}{c1:fig:entity}

\code{std_logic} är signaltypen för en enstaka bit. Den är inte inbyggd i VHDL utan kommer från
paketet \code{std_logic_1164}, vilket är därför varje fil i den här boken inleds med:

\begin{vhdlcode}
library ieee;
use ieee.std_logic_1164.all;
\end{vhdlcode}

Den har nio värden i stället för två, eftersom verklig hårdvara behöver mer än sant och falskt:
\begin{itemize}
  \item \code{'0'} och \code{'1'}: driven låg och driven hög. De enda två du kommer att skriva i
    den här boken.
  \item \code{'Z'}: högimpedans, ingenting driver ledningen alls.
  \item \code{'U'}: oinitierad, det en signal håller innan något har tilldelat den ett värde.
  \item \code{'X'}: okänd, det en simulator visar när två saker driver samma signal och är oense.
  \item \code{'-'}, \code{'W'}, \code{'L'}, \code{'H'}: don't-care-värden och svagt drivande
    värden, sällan skrivna för hand.
\end{itemize}

Att känna igen de sju du inte kommer att skriva spelar ändå roll: de är hur en simulator talar om
för dig att något är fel. En signal som ligger kvar på \code{'U'} eller \code{'X'} är en bugg,
inte ett värde.

Arkitekturen fyller i lådan:

\begin{vhdlcode}
architecture behaviour of or_gate is
begin
    x <= a or b;
end architecture;
\end{vhdlcode}

\lecfig[0.62]{lectures/L01/appendix/images/or_gate_arch.png}{Arkitekturen \code{or_gate}:
implementationen bakom entiteten.}{c1:fig:arch}

Tillsammans utgör de den kompletta modulen:

\lecfig[0.62]{lectures/L01/appendix/images/or_gate_module.png}{Modulen \code{or_gate}: entitet och
arkitektur tillsammans.}{c1:fig:module}

\code{x <= a or b;} är \secref{c1:sec:boolean}:s booleska algebra nästan ordagrant: VHDL:s
\code{and}, \code{or}, \code{not} och \code{xor} arbetar direkt på \code{std_logic}, så varje
ekvation du härlett för hand blir VHDL med knappt någon översättning alls.

Tilldelningen är \textbf{konkurrent}. Den körs inte en gång; den gäller kontinuerligt, som en
ledning mellan verkliga grindar.

Nyckelordet \code{signal} deklarerar en intern ledning inuti en arkitektur, till skillnad från en
utifrån synlig port. Du använder en i Övning~\ref{c1:ex:adas} för att namnge ett mellanresultat i
stället för att nästla in det på plats, och \chapref{c2:ch} lutar sig mot samma idé rakt igenom.

% ------------------------------------------------------------------------------------------
\section{Den kompletta modulen: \code{or_gate}}
\label{c1:sec:module}

\Secref{c1:sec:vhdl}:s delar tillsammans blir en komplett, syntetiserbar fil. Det här är hela
\file{or_gate.vhd}, och det är en riktig konstruktion som syntetiseras och körs på DE0-CV:

\vhdlfile{lectures/L01/or_gate/or_gate.vhd}

\begin{itemize}
  \item \code{library}/\code{use}-satserna drar in \code{std_logic}.
  \item \code{entity} deklarerar utsidan: två ingångar, en utgång.
  \item \code{architecture} implementerar insidan, som en enda konkurrent tilldelning.
\end{itemize}

Det är den minsta kompletta VHDL-konstruktion som finns, och formen ändras aldrig. Varje modul i
den här boken, ända upp till CAN-kontrollern i \chapref{c17:ch}, är samma
\code{library}/\code{entity}/\code{architecture}-skelett. Det som växer är arkitekturkroppen,
aldrig ramen.

\paragraph{Att kontrollera ditt arbete}
\code{or_gate} har en utdelad \file{or_gate_tb.vhd}, som driver alla fyra ingångskombinationerna
och kontrollerar utgången:

\begin{shell}
cd lectures/L01/or_gate
ghdl -a --std=93 or_gate.vhd or_gate_tb.vhd
ghdl -e --std=93 or_gate_tb
ghdl -r --std=93 or_gate_tb --assert-level=error --stop-time=10ms
\end{shell}

Ingen utskrift utöver en avslutande notering betyder att varje kontroll gick igenom. Testbänkar
behandlas ordentligt i \secref{c2:sec:testbenches}; tills vidare, betrakta de här tre kommandona
som sättet du bekräftar att en modul fungerar på.

% ------------------------------------------------------------------------------------------
\section{En bromsassistent, från krav till krets}
\label{c1:sec:adas}

Passet bygger en annan krets än \code{or_gate}, live, så att du får se samma väg gås två gånger:
en gång här i skrift, och en gång på något du inte redan läst svaret på.

Kravet: \textbf{bilen bromsar om föraren bromsar, eller om assistanssystemet bestämmer sig för
det.} Assistanssystemet bestämmer sig för det när endera detektorn ser något, om det inte
samtidigt rapporterar ett fel, i vilket fall det inte anförtros att bromsa alls.

\begin{booktable}{@{}p{8em}p{4em}L@{}}
  \thead{Port} & \thead{Riktning} & \thead{Betydelse} \\
  \midrule
  \code{driver_brake} & in & Föraren står på bromspedalen. \\
  \code{sensor} & in & Avståndssensorn ser ett hinder. \\
  \code{radar} & in & Radarn ser ett hinder. \\
  \code{error} & in & Assistanssystemet rapporterar ett fel. \\
  \code{engine_brake} & out & Lägg an bromsarna. \\
\end{booktable}

Tre satser avgör tillsammans hela strukturen:
\begin{itemize}
  \item endera detektorn räcker på egen hand, och ingen av dem väger tyngre än den andra.
  \item ett fel utlöser inget larm; det \emph{hindrar} assistanssystemet från att kunna bromsa.
  \item föraren kan alltid bromsa, även när assistanssystemet är trasigt och inte gör någonting.
\end{itemize}

Den sista är ett säkerhetskrav, och det är den att hålla fast vid när du härleder nätet: pedalen
är inte en ingång till assistanslogiken, den går förbi den. Räkna ut vad som ändras om pedalen i
stället matas genom fellogiken, så har du hittat buggen som strukturen finns till för att
förhindra.

\subsection*{Sanningstabellen}
Sexton rader, i den ordning en fyrabitars räknare skulle producera dem. \code{adas_brake} är ingen
port; den är assistanssystemets eget beslut, visad som en mellankolumn så att du kan se var
strukturen kommer ifrån.

\begin{narrowtable}{cccccc}
  \thead{\code{driver_brake}} & \thead{\code{sensor}} & \thead{\code{radar}} &
    \thead{\code{error}} & \thead{\code{adas_brake}} & \thead{\code{engine_brake}} \\
  \midrule
  0 & 0 & 0 & 0 & 0 & 0 \\
  0 & 0 & 0 & 1 & 0 & 0 \\
  0 & 0 & 1 & 0 & 1 & 1 \\
  0 & 0 & 1 & 1 & 0 & 0 \\
  0 & 1 & 0 & 0 & 1 & 1 \\
  0 & 1 & 0 & 1 & 0 & 0 \\
  0 & 1 & 1 & 0 & 1 & 1 \\
  0 & 1 & 1 & 1 & 0 & 0 \\
  1 & 0 & 0 & 0 & 0 & 1 \\
  1 & 0 & 0 & 1 & 0 & 1 \\
  1 & 0 & 1 & 0 & 1 & 1 \\
  1 & 0 & 1 & 1 & 0 & 1 \\
  1 & 1 & 0 & 0 & 1 & 1 \\
  1 & 1 & 0 & 1 & 0 & 1 \\
  1 & 1 & 1 & 0 & 1 & 1 \\
  1 & 1 & 1 & 1 & 0 & 1 \\
\end{narrowtable}

Den nedre halvan är genomgående \code{1}: så snart \code{driver_brake} är satt ändrar ingenting
annat svaret. Det är hur ”åsidosätter allt” ser ut i en tabell.

\subsection*{Fyra svar att ta med sig till passet}
Allt ovan är specifikationen. Härledningen, grindnätet och VHDL-koden är det som byggs live, och
det är värt att komma dit med ett eget svar att jämföra mot snarare än ett du redan läst. Så,
utifrån tabellen ovan:
\begin{itemize}
  \item Läs av en summa-av-produkter-ekvation för \code{engine_brake} direkt ur den, på samma sätt
    som \secref{c1:sec:boolean} gjorde.
  \item Förenkla den för hand. Kravet är tre satser långt, så det minimala nätet är avsevärt
    mindre än vad den radvisa avläsningen ger dig. Räkna ut hur mycket mindre innan du får det
    berättat för dig.
  \item Räkna grindarna din förenklade ekvation behöver, och skriv ner antalet.
  \item Skissa nätet, och avgör vilken enda grind säkerhetsargumentet ovan vilar på.
\end{itemize}

Ta med de fyra svaren till passet. Att ha fel om något av dem är värt mer än att inte ha gissat,
eftersom du då vet exakt vilket steg du ska granska om när den live-gjorda härledningen når fram
till det. Det är \secref{c1:sec:cv}:s ”härled först, simulera sedan” tillämpat på en krets du inte
byggt ännu, och det är den vana resten av boken är byggd på. Efteråt skriver du VHDL-koden själv,
i Övning~\ref{c1:ex:adas}, och en testbänk avgör om dina ekvationer och sanningstabellen ovan är
överens.

\subsection*{Varför det finns en testbänk för den}
Sexton kombinationer är få nog att kontrollera för hand en gång. Det är inte få nog att
kontrollera varje gång någon redigerar filen, vilket är den situation varje modul i den här boken
befinner sig i härifrån och framåt.

En \textbf{testbänk} är en andra VHDL-fil som gör kontrollen. Den är inte en del av
konstruktionen och når aldrig FPGA:n: den driver ingångarna genom varje fall, räknar ut vad
utgången borde ha varit, jämför och rapporterar. För \code{adas}, som skrivs vid sidan av den
under passet, är det en loop över de sexton kombinationerna som återger bromsregeln oberoende.

\emph{Oberoende} är det viktiga ordet. En testbänk som läste sitt förväntade svar ur
konstruktionen skulle stämma med den per konstruktion, även när konstruktionen har fel. Att
återge kravet separat är det som gör det till en kontroll.

Det är därför nästan varje katalog i kursrepot har en \file{_tb.vhd} vid sidan av modulen: det är
så du kan få en modul i handen, ändra den, och inom en sekund veta om du gjort sönder den.
\Secref{c2:sec:testbenches} går igenom hur de körs; du ombeds inte skriva någon någonstans i den
här boken.
